% Source-preserving cumulative Indonesian driver for Fremlin Measure Theory 111-114.
% The build staging directory supplies frozen Volume 1 dependencies and overlays
% source/id-ID/mt111.tex through source/id-ID/mt114.tex.
\input volwp.aux
\mtlulustyle
\luluvolumeno=1

% Localize reader-facing headings generated by frozen source macros.
\def\Notesheader#1{\immediate\write0{#1 Catatan dan komentar}
     \wheader{#1 Catatan}{16}{8}{5}{48pt}
     \noindent{\bf #1 Catatan dan komentar}}
\def\vNotesheader#1#2{\immediate\write0{#2 Catatan dan komentar}
     \wheader{#2 Catatan}{16}{8}{5}{#1}
     \noindent{\bf #2 Catatan dan komentar}}
\def\Hint#1{({\it Petunjuk\/}: #1)}
\long\def\proof#1{\ifwithproofs{\medskip

     \noindent{\bf Bukti} #1}\else\fi}

% Preserve the four source diagrams used by Section 113.
\input mt113-dvipdfmx-images

\special{pdf:docinfo <<
  /Title (Fondasi Teori Ukuran - Volume 1, Bagian 111-114)
  /Author (D. H. Fremlin; adaptasi Bahasa Indonesia atas arahan pengguna)
  /Subject (Adaptasi Bahasa Indonesia dari Measure Theory, Volume 1, Bagian 111-114)
  /Keywords (teori ukuran, aljabar sigma, ruang ukur, ukuran luar, ukuran Lebesgue, id-ID, O007)
>>}
\begingroup
\catcode64=12
\special{pdf:put @catalog << /Lang (id-ID) /ViewerPreferences << /DisplayDocTitle true >> >>}
\endgroup

\pageno=1
\vbox to \vsize{
  \vskip 42pt
  \centerline{\fourteenbf FONDASI TEORI UKURAN}
  \vskip 15pt
  \centerline{\bf Volume 1: Minimum yang Tak Tereduksi}
  \vskip 42pt
  \centerline{\fourteenbf Bagian 111-114}
  \vskip 8pt
  \centerline{\fourteenbf Dari Aljabar sigma menuju Ukuran Lebesgue}
  \vskip 48pt
  \centerline{D. H. Fremlin}
  \vskip 6pt
  \centerline{Adaptasi Bahasa Indonesia atas arahan pengguna}
  \vfill
  \hsize=330pt
  \leftskip=45pt
  \rightskip=45pt
  \noindent Sumber: \it Measure Theory, Volume 1: The Irreducible Minimum,
  \rm Bagian 111-114, D. H. Fremlin. Terjemahan dan modernisasi pembaca:
  22 Agustus 2026. Materi turunan Fremlin tetap berada di bawah Design
  Science License. Sumber yang dapat diedit, lisensi, atribusi, pemetaan
  segmen, dan catatan perubahan disertakan bersama edisi ini.
  \vskip 18pt
}
\eject

\pageno=1
\def\folio{\number\pageno}
\headline={\rlap{\eightbf 111}\hfill{\eightit Aljabar-$\sigma$}\hfill}
\footline={\hfill\eightrm\folio\hfill}
\input mt111
\eject
\headline={\rlap{\eightbf 112}\hfill{\eightit Ruang ukur}\hfill}
\input mt112
\eject
\headline={\rlap{\eightbf 113}\hfill{\eightit Ukuran luar dan konstruksi Carath\'eodory}\hfill}
\input mt113
\eject
\headline={\rlap{\eightbf 114}\hfill{\eightit Ukuran Lebesgue pada $\Bbb R$}\hfill}
\input mt114
\end
